% Volume III, Chapter 6 -- addendum: localization over the non-analytic group ring Z_2[[P]]
% (open problem (a) of Section~\ref{V3P6:sec:retro}).  Verified by code/v3p6_pro2_localization.py
% and, for the 2-adic orders, code/v3p6_level5_v2.py.
\section{Localization over the non-analytic group \texorpdfstring{$\Z_2[[P]]$}{Z2[[P]]}}\label{V3P6:sec:pro2}

Corollary~\ref{V3P6:cor:status} leaves parts (i$'$) and (ii) of Theorem~\ref{V3P6:thm:main} open
for the canonical tower, because nothing is known about $\Lambda=\Z_2[[P]]$ beyond what holds for
every pro-$2$ group. This section separates what can be proved about $\Lambda$ from what cannot,
and proves an unconditional substitute for (ii): a compatible system of characteristic elements
at the finite levels, whose boundaries are the finite parts of the Iwasawa module and whose
reduced norms are the Artin factors of the class number formula. Throughout, $G_m=P/P_m$,
$n_m=3|G_m|$, $N_m=\sum_{g\in G_m}g\in\Z_2[G_m]$, and $\DeltaH_\Lambda\in M_3(\Lambda)$ is the
equivariant operator of Section~\ref{V3P6:sec:main} for the tower $Y_m\to X_3$, so that
$\Z_2[G_m]\otimes_\Lambda\DeltaH_\Lambda=\DeltaH(Y_m)\otimes\Z_2$ under
$\Z_2[G_m]^{3}=\Z_2^{V(Y_m)}$.

\subsection*{What is known about \texorpdfstring{$\Lambda$}{Lambda}}

\begin{proposition}[The group $P$ and its ring]\label{V3P6:prop:pro2basic}
\begin{enumerate}
\item[(i)] $P$ is a torsion-free pro-$2$ group, topologically generated by $3$ elements, of
infinite rank; $d(P_k)\ge2k$.
\item[(ii)] $\Lambda=\varprojlim_m\Z_2[G_m]$ is a compact local ring with maximal ideal
$(2,I_P)$, $I_P$ the augmentation ideal, and every finitely generated compact
$\Lambda$-module $M$ satisfies $M=\varprojlim_m\Z_2[G_m]\otimes_\Lambda M$.
\item[(iii)] The necessary condition for $\Lambda$ to be a domain holds: $\Z_2[[P]]$ has no
zero divisor of the form $(g-1)\nu$ with $g$ of finite order and $\nu$ its norm element.
\end{enumerate}
\end{proposition}

\begin{proof}
(i) Generation, rank and $d(P_k)$ are Theorem~\ref{V3P6:thm:tower}. A pro-$2$ group has no
torsion of odd order, so let $x\in1+\varphi\mathcal O$ with $x^{2^{k}}\in K^{\times}$,
$K=\F_2((t))$ the centre of $D_0$. The field $K(x)\subseteq D_0$ has degree $1$ or $3$ over $K$,
since $D_0$ is a division algebra of degree $3$. Over $K(x)$, in characteristic $2$, $y^{2^{k}}-x^{2^{k}}=(y-x)^{2^{k}}$, and the minimal
polynomial $f$ of $x$ over $K$ divides it, so $f=(y-x)^{j}$ with $j=[K(x):K]\in\{1,3\}$. If $j=3$,
then $f=y^{3}+xy^{2}+x^{2}y+x^{3}$ has its coefficients in $K$, so $x\in K$, a contradiction.
Hence $j=1$, $x\in K^{\times}\cap(1+\varphi\mathcal O)=1+t\F_2[[t]]$, and $x$ is
trivial in $P$. (ii) is standard for pro-$p$ groups: $\Lambda/(2,I_P)=\F_2$, and every element
outside $(2,I_P)$ is a unit, being a unit in every $\Z_2[G_m]$, which is local because $G_m$ is
a $2$-group; the limit statement holds because $M$ is a quotient of $\Lambda^{r}$ and inverse
limits of compact modules are exact. (iii) If $g$ has order $p$, then
$(g-1)(1+g+\dots+g^{p-1})=0$; by (i) there is no such $g\ne1$.
\end{proof}

The ring-theoretic questions that (i$'$) and (ii) need are open for $\Lambda$.

\begin{remark}[What is not known]\label{V3P6:rem:pro2open}
\begin{enumerate}
\item[(a)] \emph{Domain.} Neumann's theorem \cite{Neu88} says that $\Z_p[[G]]$ is a domain when
$G$ is a torsion-free \emph{$p$-adic analytic} pro-$p$ group; its proof goes through Lazard's
filtration and the fact that the associated graded ring is a polynomial ring \cite{Laz65}. $P$ is
not analytic (Theorem~\ref{V3P6:thm:tower}(iii)), so the theorem does not apply.
Proposition~\ref{V3P6:prop:pro2basic}(i) supplies the necessary condition, torsion-freeness, but
we know of no theorem that makes $\Z_2[[P]]$ a domain. Lazard's argument would go through if the
graded ring of $\F_2[[P]]$ for the filtration induced by the $P_m$ were a domain. The $P_m$ form a
restricted $N$-series, so by Jennings and Quillen this graded ring is the restricted enveloping
algebra of the graded restricted Lie algebra $\operatorname{gr}P$. In $\operatorname{gr}P$ the
squaring map is injective on every homogeneous piece: the square of $1+a\varphi^{k}$ is
$1+a\sigma^{k}(a)\varphi^{2k}$, and $a\sigma^{k}(a)$ is nonzero modulo the centre. So the obvious
source of zero divisors, $x^{[2]}=0$, is absent, but whether the enveloping algebra is a domain is
not known to us.
\item[(b)] \emph{Fractions.} Even for a domain, the nonzero elements need not form an Ore set.
Goldie's theorem gives the skew field $Q(\Lambda)$ for Noetherian domains, and whether
$\Z_2[[P]]$ is Noetherian is open (Corollary~\ref{V3P6:cor:status}); to our knowledge no pro-$p$ group of
infinite rank is known to have a Noetherian completed group ring. The Cohn localization at $\Sigma=\Lambda\setminus\{0\}$
exists by its universal property, but the localization sequence of
Definition~\ref{V3P6:def:cohn} needs $\Lambda\to\Sigma^{-1}\Lambda$ to be injective and stably flat,
which we cannot verify.
\item[(c)] \emph{Consequence.} The element $\mathcal L=[\DeltaH_\Lambda]\in K_1(Q(\Lambda))$ and the
identity $\partial\mathcal L=[X(Y/P)]$ are available only conditionally. Theorem~\ref{V3P6:thm:main}(i)
does hold unconditionally: $\DeltaH_\Lambda$ is injective on $\Lambda^{3}$.
\end{enumerate}
\end{remark}

\subsection*{An unconditional substitute: the tower of finite characteristic elements}

The trivial representation is the only obstruction at a finite level, so we remove it. Put
\[
A_m=\Z_2[G_m]/\Z_2N_m,\qquad \alpha_m=\text{image of }\DeltaH_\Lambda\text{ in }M_3(A_m),
\qquad X_m^{\mathrm{nt}}=\coker\bigl(\alpha_m\colon A_m^{3}\to A_m^{3}\bigr).
\]
$A_m$ is a ring, because $\Z_2N_m$ is a two-sided ideal ($gN_m=N_mg=N_m$). It is a $\Z_2$-order
in the semisimple algebra $\Q_2A_m=(1-e_m)\Q_2[G_m]=\bigoplus_{\rho\ne1}$, where
$e_m=N_m/|G_m|$ is the trivial idempotent; in particular $A_m$ is Noetherian, free of rank
$|G_m|-1$ over $\Z_2$, and $\Q_2A_m=A_m[\tfrac12]$ is its localization at the central
multiplicative set $\{2^{k}\}$. For such an order the localization sequence is classical and
unconditional:
\[
K_1(A_m)\to K_1(\Q_2A_m)\xrightarrow{\ \partial\ }K_0(A_m,\Q_2)\to K_0(A_m)\to K_0(\Q_2A_m),
\]
where $K_0(A_m,\Q_2)$ is the Grothendieck group of finite $A_m$-modules of projective dimension
$\le1$ and $\partial[\alpha]=[\coker\alpha]$ for injective $\alpha\in M_r(A_m)$, with the sign
convention of Definition~\ref{V3P6:def:cohn}.

\begin{lemma}[The rings $A_m$ form the tower]\label{V3P6:lem:pro2limit}
The projections $G_{m+1}\to G_m$ induce surjective ring maps $\pi_m\colon A_{m+1}\to A_m$ with
$\pi_m(\alpha_{m+1})=\alpha_m$, and $\Lambda\to\varprojlim_mA_m$ is an isomorphism of
topological rings.
\end{lemma}

\begin{proof}
$N_{m+1}\mapsto|P_m/P_{m+1}|\,N_m=2^{d_m}N_m$, so the ideal $\Z_2N_{m+1}$ maps into $\Z_2N_m$,
and $\pi_m$ is well defined and surjective. Both $\alpha_{m+1}$ and $\alpha_m$ are images of
$\DeltaH_\Lambda$. If $\lambda\in\Lambda$ maps to $0$ in every $A_m$, its images in
$\Z_2[G_m]$ are $c_mN_m$ with $c_m=2^{d_m}c_{m+1}$, so $c_m$ is divisible by every power of $2$,
and $\lambda=0$. Surjectivity follows from compactness, since all transition maps are
surjective.
\end{proof}

\begin{theorem}[Finite characteristic elements along the canonical tower]\label{V3P6:thm:pro2}
For every $m\ge1$, unconditionally:
\begin{enumerate}
\item[(i)] $\alpha_m$ is invertible in $M_3(\Q_2A_m)$, the sequence
$0\to A_m^{3}\xrightarrow{\alpha_m}A_m^{3}\to X_m^{\mathrm{nt}}\to0$ is exact, and
\[
X_m^{\mathrm{nt}}=\coker\bigl[\DeltaH(Y_m)\otimes\Z_2\ \big|\ \mathbf 1_{V_0},\mathbf 1_{V_1},
\mathbf 1_{V_2}\bigr]
\]
is a finite $A_m$-module of projective dimension $\le1$, where $\mathbf 1_{V_s}$ is the indicator
of the vertices of type $s$.
\item[(ii)] $\mathcal L_m:=[\alpha_m]\in K_1(\Q_2A_m)$ satisfies $\partial\mathcal L_m=[X_m^{\mathrm{nt}}]$,
and any $\mathcal L'\in K_1(\Q_2A_m)$ with $\partial\mathcal L'=[X_m^{\mathrm{nt}}]$ differs from
$\mathcal L_m$ by an element of the image of $K_1(A_m)$.
\item[(iii)] (Compatibility.) $\pi_{m*}\mathcal L_{m+1}=\mathcal L_m$ and
$A_m\otimes_{A_{m+1}}X_{m+1}^{\mathrm{nt}}\cong X_m^{\mathrm{nt}}$ with
$\operatorname{Tor}_1^{A_{m+1}}(A_m,X_{m+1}^{\mathrm{nt}})=0$, so the boundaries are compatible:
$\pi_{m*}\partial\mathcal L_{m+1}=\partial\mathcal L_m$. Hence
\[
\mathcal L_\infty:=(\mathcal L_m)_m\in\varprojlim_mK_1(\Q_2A_m),\qquad
X(Y/P)=\coker\DeltaH_\Lambda\cong\varprojlim_mX_m^{\mathrm{nt}} .
\]
\item[(iv)] (Artin descent.) Over $\C$, $\Q_2A_m\otimes\C=\bigoplus_{\rho\ne1}M_{\dim\rho}(\C)$ and
the reduced norm of $\mathcal L_m$ has components $\det\DeltaH_\rho$, $\rho\ne1$. The order of the
boundary is the class number formula:
\[
\log_2|X_m^{\mathrm{nt}}|=v_2\prod_{\rho\ne1}\det(\DeltaH_\rho)^{\dim\rho}
=v_2\Bigl([P:P_m]\,\frac{|\JacH(Y_m)|}{|\JacH(X_3)|}\Bigr)
=\log_2|G_m|+v_2|\JacH(Y_m)| .
\]
\end{enumerate}
\end{theorem}

\begin{proof}
(i) $\Q_2A_m\otimes\C$ is the sum of the nontrivial isotypic components of $\C[G_m]$, and on the
$\rho$-component $\alpha_m$ acts as $\DeltaH_\rho$ (with multiplicity $\dim\rho$), which is
nonsingular by Theorem~\ref{V3P6:thm:main}(iv). So $\alpha_m$ is invertible over $\Q_2A_m$ and
injective on the lattice $A_m^{3}$, whose cokernel is then finite; the exact sequence is a
resolution of length $1$ by free modules. Under $\Z_2[G_m]^{3}=\Z_2^{V(Y_m)}$ the element $N_me_s$
is $\mathbf 1_{V_s}$, so $A_m^{3}=\Z_2^{V(Y_m)}/\langle\mathbf 1_{V_0},\mathbf 1_{V_1},\mathbf 1_{V_2}\rangle$
and $X_m^{\mathrm{nt}}=\Z_2^{V(Y_m)}/(\operatorname{im}\DeltaH(Y_m)+\langle\mathbf 1_{V_s}\rangle)$.

(ii) is the definition of $\partial$ and the exactness of the localization sequence at
$K_1(\Q_2A_m)$.

(iii) $\pi_m(\alpha_{m+1})=\alpha_m$ by Lemma~\ref{V3P6:lem:pro2limit}. Applying
$A_m\otimes_{A_{m+1}}-$ to the resolution of (i) gives $A_m^{3}\xrightarrow{\alpha_m}A_m^{3}\to
A_m\otimes X_{m+1}^{\mathrm{nt}}\to0$, and $\alpha_m$ is injective, so the Tor term vanishes and
the base change is $X_m^{\mathrm{nt}}$. The localization sequences are natural for the ring map
$\pi_m$, which sends $2$ to $2$, so $\partial$ commutes with $\pi_{m*}$ (on $K_0(-,\Q_2)$ the map is
$[M]\mapsto[A_m\otimes^{\mathbf L}M]$). For the limit, $X=\varprojlim X_m$ with
$X_m=\Z_2[G_m]\otimes_\Lambda X=\coker\DeltaH(Y_m)\otimes\Z_2$
(Proposition~\ref{V3P6:prop:pro2basic}(ii)), and $X_m\to X_m^{\mathrm{nt}}$ is surjective with kernel
$K_m$, the $\Z_2$-span of the images of the $\mathbf 1_{V_s}$. The transition maps send
$\mathbf 1_{V_s}(Y_{m+1})\mapsto2^{d_m}\mathbf 1_{V_s}(Y_m)$, so a compatible family in $(K_m)$
lies in $\bigcap_j2^{j}K_m=0$ (Krull, $K_m$ being a finitely generated $\Z_2$-module); hence
$\varprojlim K_m=0$, and $\varprojlim X_m\cong\varprojlim X_m^{\mathrm{nt}}$ because $\varprojlim$
is exact on compact modules.

(iv) The first statement is the decomposition used in (i). For the order: $A_m^{3}$ is a free
$\Z_2$-module of rank $3(|G_m|-1)$, and for an injective endomorphism of a lattice
$|\coker\alpha_m|=|\det_{\Q_2}\alpha_m|_2^{-1}$. The $\Q_2$-determinant of $\alpha_m$ on
$\Q_2A_m^{3}$ is the determinant of $\DeltaH(Y_m)$ on the nontrivial isotypic part, i.e.\
$\prod_{\rho\ne1}\det(\DeltaH_\rho)^{\dim\rho}$, which equals $[P:P_m]|\JacH(Y_m)|/|\JacH(X_3)|$
by Theorem~\ref{V3P6:thm:main}(iv). Finally $|\JacH(X_3)|=147$ is odd.
\end{proof}

\begin{remark}[The orders, and the checks]\label{V3P6:rem:pro2check}
With Theorems~\ref{V3P6:thm:v2five} and~\ref{V3P6:thm:v2eight}, (iv) gives
$\log_2|X_m^{\mathrm{nt}}|=3+18=21$, $6+99=105$, $8+289=297$, $11+1742=1753$, $14+11595=11609$,
$16+36937=36953$ and $19+254662=254681$ for $m=2,\dots,8$. Over $\Z$ the same construction gives
$|X_2^{\mathrm{nt}}|=8\cdot|\JacH(Y_{24})|/147=2^{21}7^{14}$, which is the product of the seven
twisted determinants of the layer $Y_2\to Y_1$. The script \texttt{code/v3p6\_pro2\_localization.py}
(numpy; about three minutes) checks the ingredients at the levels $m\le5$: the orders $|G_m|$ and
the central elementary abelian layers (OR); the push-forward and pull-back identities
$\Pi\,\DeltaH(Y_{m+1})=\DeltaH(Y_m)\,\Pi$ and $\DeltaH(Y_{m+1})\Pi^{t}=\Pi^{t}\DeltaH(Y_m)$ behind
Lemma~\ref{V3P6:lem:pro2limit} (PF); the block decomposition of each layer (BL); the sector lemma
for all twisted blocks (SE); the Smith forms $\coker\DeltaH(Y_m)=\Z\oplus\JacH(Y_m)$ and the order
$|X_m^{\mathrm{nt}}|=|G_m||\JacH(Y_m)|/147$ over $\Z$ for $m=2,3$ (SM); the relative class number
formula of the layers $Y_2/Y_1$, $Y_3/Y_2$ and $Y_4/Y_3$ with exact integers, $|\JacH(Y_4)|$ being
reconstructed by Chinese remaindering (CN); and the surjectivity of
$\JacH(Y_{m+1})\to\JacH(Y_m)$ with kernel of order $\prod_{\chi\ne1}\det\DeltaH_\chi/|H_m|$ (KR).
\end{remark}

\begin{remark}[How far this goes]\label{V3P6:rem:pro2scope}
Theorem~\ref{V3P6:thm:pro2} is the form of Theorem~\ref{V3P6:thm:main}(i$'$),(ii) that holds without
knowing anything about $\Z_2[[P]]$: the characteristic element lives in the inverse limit
$\varprojlim K_1(\Q_2A_m)$ instead of $K_1(Q(\Lambda))$, its boundary is the system of finite
modules $X_m^{\mathrm{nt}}$ whose limit is $X(Y/P)$, and its reduced norms are the Artin factors. This
is the shape of the finite-level formulations of non-commutative Iwasawa theory, where $K_1$ of a
localized Iwasawa algebra is approached through the finite group rings. If $\Lambda$ turns out to be
an Ore domain, then $\mathcal L=[\DeltaH_\Lambda]\in K_1(Q(\Lambda))$ exists and
$\partial\mathcal L=[X]$ by Theorem~\ref{V3P6:thm:main}(ii); we do not claim a map relating
$\mathcal L$ to $\mathcal L_\infty$, because a skew field of fractions of $\Lambda$ need not map to
the finite-level algebras $\Q_2A_m$. Open
problem (a) of Section~\ref{V3P6:sec:retro} is therefore sharpened to: is $\Z_2[[P]]$ a Noetherian
domain, or at least an Ore domain?
\end{remark}

\begin{remark}[Corrections to the proposed programme]\label{V3P6:rem:pro2corrections}
(1) Neumann's theorem does not say that $\Z_2[[P]]$ is a domain: it requires a $p$-adic analytic
group, and $P$ is not one. What is proved is that $P$ is torsion-free
(Proposition~\ref{V3P6:prop:pro2basic}), which is necessary, not sufficient. (2) The Ore
localization at $\Lambda\setminus\{0\}$, i.e.\ the skew field $Q(\Lambda)$, is not known to exist:
Goldie's theorem needs a Noetherian ring (Remark~\ref{V3P6:rem:pro2open}). (3) Therefore
$\partial[\DeltaH_\Lambda]=[X(Y/P)]$ in $K_1(Q(\Lambda))$ is not proved unconditionally. What holds
unconditionally is Theorem~\ref{V3P6:thm:pro2}, the same identity at every finite level, compatibly
along the tower, with $X(Y/P)$ recovered as the limit. (4) The descent formula with denominator
$|\JacH(X_3)|$ is Theorem~\ref{V3P6:thm:main}(iv); in Theorem~\ref{V3P6:thm:pro2}(iv) it computes the
order of the boundary. (5) $d(P_k)\ge2k$ and the absence of the Ore set of \cite{CFKSV} are
Theorems~\ref{V3P6:thm:tower}(iii) and~\ref{V3P6:thm:nozp}, as stated in the proposed programme.
\end{remark}
